\section{The 4-cycle and other graphs}\label{sec:examples}

In this section $Q$ is symmetric. So the admissible faces are the connected
faces with positive start mass, and they satisfy (H)
(Lemma~\ref{lem:admissible}). By Theorem~\ref{thm:hull}(c),
Theorem~\ref{thm:window} applies at every tie. We call the faces in
$\mathcal F_\tau$ the \emph{winners} at $\tau$. We first complete the worked
example of Section~\ref{sec:worked} with all its constants, then show on $K_3$
that the start can decide the winner, and then list the outcomes for other
graphs.

\begin{corollary}[the 4-cycle cascade]\label{cor:C4}
Let $Q$ be the 4-cycle $C_4$ with $q_{i,i\pm1}=1$ (indices mod $4$), and let
$\mu$ be uniform.
\begin{enumerate}[(a)]
\item The admissible faces are the vertices, the adjacent pairs, the arcs of
three states and $C_4$, with exit rates $2$, $1$, $2-\sqrt2$ and $0$. The
winners are the vertices for $\tau<1$, the adjacent pairs for $1<\tau<2$ and
$C_4$ for $\tau>2$. The arcs of three states never win.
\item $K_{\rm vertex}=\frac14$, $K_{\rm pair}=1/\sqrt{2\pi}$,
$K_{\rm arc3}=(4+3\sqrt2)/(4\pi)$ and $K_{C_4}=8/\pi^{3/2}$, and
$\det\Sigma=1/512$ on $C_4$.
\item For every $N\ge2$ and $0<h<\frac12$, $M_{\rm pair}=M_{\rm vertex}S_N(u)$
with $u=h/(1-2h)$. So the face maxima of a vertex and an adjacent pair cross
exactly at $u=u_N$, the crossing of~\cite{paperA}. For large $N$ this
crossing satisfies $T+\frac12\log T=L+c_0+O(1/L)$.
\item For large $N$, every crossing $T_N\in[1,N^{1/3}]$ of the face maxima of
an adjacent pair and $C_4$ satisfies $T_N+\log T_N=2L-\log(8\sqrt2/\pi)+O(1/L)$.
\item Put $b=-\frac12\log2-\frac12\log(8\sqrt2/\pi)$. In the tie window at
$\tau=1$ the face that carries the global mode changes from the vertices to
the adjacent pairs only at $T=L-\frac12\log L+c_0+o(1)$. In the tie window at
$\tau=2$ it changes from the adjacent pairs to $C_4$ only at
$T=2(L-\frac12\log L+b)+o(1)$.
\end{enumerate}
\end{corollary}

In words, the mode moves from a vertex to an adjacent pair at the crossing of
Kagey's board, exactly at every $N$, and from the pairs to the whole cycle at a
crossing whose equation has the coefficient $1=\frac\Delta2$ of $\log T$, since
the cycle has two more states than a pair.

\begin{proof}
(a) A proper face that is not an arc is disconnected, so it is not
admissible. For an adjacent pair $-Q_F$ has rows $(2,-1)$ and $(-1,2)$, with
least eigenvalue $1$. For an arc of three states $-Q_F=2I-A$, where $A$ is the
adjacency matrix of the path on three states, and $r=(1/\sqrt2,1,1/\sqrt2)$ is a
positive eigenvector with eigenvalue $2-\sqrt2$. Also $Q\ones=0$. So
$\Phi_\tau$ equals $2\tau$, $\tau+1$, $(2-\sqrt2)\tau+2$ and $3$, and comparing
the first, second and fourth gives the winners. Since $2-\sqrt2>\frac12$, the
3-arc has $\Phi_\tau>\frac12(\tau+1)+\frac12\cdot3$, which is at least the
smaller of the exponents of the pair and of $C_4$. So it never wins.

(b) A vertex has $K=\mu_i=\frac14$. The pair and $C_4$ have symmetric $Q_F$
with constant row sums, so Corollary~\ref{cor:symface} applies, with
$\mu(F)=\frac12$ and one spanning tree for the pair, and $\mu(F)=1$ and four
spanning trees for $C_4$. On $C_4$ the law $\alpha^*$ is uniform and every edge
has conductance $\frac14$, so Proposition~\ref{prop:constant}(b) gives
$\det\Sigma=2^3\cdot4^{-8}/(4\cdot4^{-3})=1/512$. The 3-arc has unequal row
sums, and we use Proposition~\ref{prop:constant}(b) with $\nu=\ones$. Here
$\alpha^*=r\circ r/\lvert r\rvert^2=(\frac14,\frac12,\frac14)$ and
$r_\alpha=l_\alpha=\sqrt{\alpha^*}$, so
$(\mu\cdot r_\alpha)(l_\alpha\cdot\ones)=\frac14(1+1/\sqrt2)^2=(3+2\sqrt2)/8$.
The only spanning tree is the path, with $\tree(c)=c_{12}c_{23}=\frac18$, and
$\prod_i\alpha^*_i=\frac1{32}$. So
$K_{\rm arc3}=\frac{3+2\sqrt2}8\cdot\frac{32}{2\sqrt2}\cdot\frac1{4\pi}
=\frac{4+3\sqrt2}{4\pi}$.

(c) The identity is Proposition~\ref{prop:pair} with $a=1$ and $q=2$. So the
two face maxima are equal only at $u=u_N$, and by
Theorem~\ref{thm:crossing}(i) this crossing lies in $[L/2,2L]$ for large $N$. As
$K_{\rm pair}/K_{\rm vertex}=2\sqrt{2/\pi}=e^{-c_0}$,
Theorem~\ref{thm:crossing}(ii) with $\Delta=\delta=1$ gives the equation.

(d) This is Theorem~\ref{thm:crossing} with $F=C_4$, $G$ an adjacent pair,
$\Delta=2$, $\delta=1$ and $K_F/K_G=8\sqrt{2\pi}/\pi^{3/2}=8\sqrt2/\pi$.

(e) At $\tau=1$ the tied faces are the vertices and the adjacent pairs, and
at $\tau=2$ they are the adjacent pairs and $C_4$. Faces of one type have the
same line. By Remark~\ref{rem:half} the kink of $E$ is at $t=b_{FG}$, which is
$c_0$ at $\tau=1$ by (c) and $b$ at $\tau=2$ by (d). So
Theorem~\ref{thm:window} gives (e).
\end{proof}

The finite parts of (a) are checked in Lean. Lean also checks that a closed
form for the constants of all arcs, which we do not use, gives the values of
$K_{\rm pair}$ and $K_{\rm arc3}$ above.

Proposition~\ref{prop:nextterm} gives the next term $c_F/T$ of the face
maxima, with the explicit coefficient
$c_F=c(\alpha^*_F)+\frac12\gamma_F^\top\Sigma_F\gamma_F$, where $c$ is given
by~\eqref{eq:c-alpha}. For an adjacent pair, Proposition~\ref{prop:nextterm}(iv)
gives $c_{\rm pair}=-\frac18$. For the 3-arc and for $C_4$ we evaluated the
coefficient exactly by computer algebra (the script
\texttt{verify\_second\_order.py}), and the values are
$c_{\rm arc3}=\frac54\sqrt2-2$ and $c_{C_4}=-\frac9{16}$. So these two
constants are exact evaluations of the formula of
Proposition~\ref{prop:nextterm}, done by computer and not by hand. With the
evaluated constant $c_{C_4}$, Proposition~\ref{prop:nextterm}(iii) gives the
next term at the crossings in (d), namely
\[
 T_N\Bigl(2L-\log T_N-T_N-\log\frac{8\sqrt2}\pi\Bigr)
 =c_{C_4}-c_{\rm pair}+O\Bigl(\frac1L\Bigr)=-\frac7{16}+O\Bigl(\frac1L\Bigr).
\]

Theorem~\ref{thm:crossing} also applies to a face that never carries the mode,
and this is a good test of the implicit equation~\eqref{eq:crossing}. Take the
3-arc $F$ against a vertex $G$, so that $\Delta=2$, $\delta=\sqrt2$ and
$\log(K_F/K_G)=\log((4+3\sqrt2)/\pi)=0.964591$. At $N=10^{14}$ a numerical
computation of the face maxima, over the symmetric compositions $(a,N-2a,a)$,
puts the crossing at $T_N=42.2633$, where
$\Delta L-\frac\Delta2\log T_N-\delta T_N=0.959070$. The root of the implicit
equation is within $0.004$ of $T_N$, while the explicit form in
Theorem~\ref{thm:crossing} gives $42.2059$. Also
$T_N(0.959070-0.964591)=-0.2333$, close to the limit
$c_{\rm arc3}-c_{\rm vertex}=-0.2322$ that
Proposition~\ref{prop:nextterm}(iii) gives with the evaluated constant
$c_{\rm arc3}$ (a vertex has $c_{\rm vertex}=0$).

By Corollary~\ref{cor:symface} the start enters $K_F$ only through $\mu(F)$, and this can
change which face wins inside a tie window.

\begin{corollary}[the start on $K_3$]\label{cor:K3start}
Let $d=3$ and $Q=J-3I$, put $T=L-\frac12\log L+t$, and let
$w=\log\frac{16}{9\sqrt3}=0.026$.
\begin{enumerate}[(a)]
\item Let $\mu=(\frac12,\frac12,0)$. For every $\eta>0$ and large $N$, if
$c_0+\eta\le t\le c_0+w-\eta$, then every global mode lies on the face $\{1,2\}$.
\item Let $\mu$ be uniform. For every compact set of $t$ and large $N$, no global mode lies on a
face of two states.
\end{enumerate}
\end{corollary}

\begin{proof}
We compare the lines of Theorem~\ref{thm:window}. All admissible faces tie at $\tau^*=1$,
because $\lambda_F=3-k_F$ (Corollary~\ref{cor:paperB}), and Theorem~\ref{thm:window} applies
since $Q$ is symmetric. By Corollary~\ref{cor:symface} the lines are
\[
\log\mu_i-2t\ \text{ for a vertex }i,\qquad
\log\bigl((\mu_i+\mu_j)\sqrt{2/\pi}\bigr)-t\ \text{ for an edge }\{i,j\},\qquad
\log K_{[3]}\ \text{ for }[3],
\]
with $K_{[3]}=3^{5/2}/(4\pi)$. A vertex line is the steepest and the line of $[3]$ is flat.

(a) The admissible faces are the vertices $1$ and $2$, the three edges and $[3]$. The edge
$\{1,2\}$ has the start mass $1$ and the other two edges have $\frac12$, so its line is the
highest edge line. It lies above the vertex line $-\log2-2t$ exactly when
$t>-\log2-\frac12\log\frac2\pi=c_0$, and above the line of $[3]$ exactly when
$t<\frac12\log\frac2\pi-\log K_{[3]}=c_0+w$. So $E$ is the line of $\{1,2\}$ on $[c_0,c_0+w]$,
its kinks are $c_0$ and $c_0+w$, and no other face has this line. The consequence in
Theorem~\ref{thm:window} gives (a).

(b) Now every vertex has the line $\log\frac13-2t$ and every edge has the line
$\log(\frac23\sqrt{2/\pi})-t$. Let $t^*$ be the point where the vertex line meets the line of
$[3]$. The gap between $\max(\log\frac13-2t,\log K_{[3]})$ and the edge line decreases for $t<t^*$
and increases for $t>t^*$, so it is least at $t^*$, and there it is
$\frac12\log(3^{7/2}/32)>0.18$. So the edge line lies below $E$ by at least $0.18$ for every $t$,
and by Theorem~\ref{thm:window}(i) the face maximum of an edge is smaller than that of a vertex
or of $[3]$, for $t$ in a compact set and large $N$.
\end{proof}

The left end $c_0$ is the crossing of Paper~A. Indeed Proposition~\ref{prop:pair} with $a=1$ and
$q=2$ shows that for every $N$ the face maxima of $\{1,2\}$ and of a vertex cross exactly at
$u=u_N$. The comparison of the lines in (a) and the value of $w$ are checked in Lean. With the
uniform start the mode jumps from the vertices straight to $[3]$, as~\cite{paperB} shows for every
complete graph. Paper~B~\cite{paperB} also compares the window in (a) with exact computations at
finite $N$, where it approaches its width $w$ slowly.

\begin{table}[t]
\centering\small
\renewcommand{\arraystretch}{1.15}
\begin{tabular}{@{}>{\raggedright\arraybackslash}p{2.1cm}
>{\raggedright\arraybackslash}p{10.2cm}
>{\raggedright\arraybackslash}p{2.6cm}@{}}
\toprule
graph & winners in order, with the value of $\tau$ where each one starts &
status\\
\midrule
$K_d$ & $\lambda_F=d-\lvert F\rvert$. Vertices for $\tau<1$ and $[d]$ for
$\tau>1$. Every face ties at $\tau=1$, and~\cite{paperB} resolves the tie at
second order. & Lean\\
symmetric $Q$ & $\lambda_{[d]\setminus\{x\}}\le q_x/(d-1)$, with equality
exactly when $q_{yx}=q_x/(d-1)$ for all $y\ne x$. With unit rates on a
connected graph, the winners go from vertices straight to $[d]$ exactly when
the graph is complete. & Lean\\
$C_d$, $d\ge4$ & Arcs of $k=1,\dots,k_d$ states with $k_d=\lfloor2d/3\rfloor$,
then $C_d$. The arc of $k\ge2$ states wins from $1/(\lambda_{k-1}-\lambda_k)$
with $\lambda_k=2-2\cos\frac\pi{k+1}$, and $C_d$ from
$\lceil d/3\rceil/\lambda_{k_d}$. These values are $1$, $1+\sqrt2$, $4.906$ and
$8.771$ for $k=2,\dots,5$, and they grow like $k^3/(2\pi^2)$. Arcs of $k\ge3$
states win exactly when $d\ge\lceil3k/2\rceil$. & Lean, asymptotics by hand\\
$P_d$, $d\ge3$ & End-arcs of $1,\dots,\lfloor2d/3\rfloor$ states, then $P_d$.
The end-arc of $k$ states has exit rate $2-2\cos\frac\pi{2k+1}$, and the first
switch is at the golden ratio $(1+\sqrt5)/2$. &
Proposition~\ref{prop:paths}\\
$P_4$ at finite $N$ & The end-arc of $3$ states never wins, and its exit rate is
only $0.0071$ too large to win. Yet with a uniform start it is the exact mode
for $T\in(9.381,12.502)$ at $N=40$, and on a window at every $N$ tested up to
$300$. & Proposition~\ref{prop:paths}, finite $N$ numerical\\
$\Theta_\varepsilon$ & Triangles $\{1,2,3\}$ and $\{4,5,6\}$ joined by a link
$3$--$4$ of rate $\varepsilon<\varepsilon_1=(64+\sqrt{46})/60$. The states
$1,2,5,6$, then $\{1,2\}$ and $\{5,6\}$ from $1$, the triangles from
$1+\varepsilon/3+O(\varepsilon^2)$, and $[6]$ from
$9/\varepsilon+2+O(\varepsilon)$. & Proposition~\ref{prop:triangles}, with 2
exact resultants\\
paw & A triangle $\{1,2,3\}$ with a leaf $4$ at $1$. $\{4\}$, then $\{1,4\}$
from $1+\sqrt2$, $\{1,2,3\}$ from $\sqrt2+\sqrt3$, and all $4$ states from
$2+\sqrt3$. & Lean\\
house & The square $1,2,3,4$ with a roof state $5$ joined to $1,2$.
$\{3\},\{4\},\{5\}$, then $\{3,4\}$ from $1$, $\{1,2,5\}$ from $1+\sqrt2$, and
all $5$ states from $2+\sqrt2$. & Lean\\
$K_4$ plus a vertex & A state $5$ joined to $1,2$ of $K_4$. $\{5\}$, then
$\{1,2,3,4\}$ from $3(\sqrt{17}+1)/8=1.921$, and all $5$ states from
$(5+\sqrt{17})/4=2.281$. & Lean\\
\bottomrule
\end{tabular}
\caption{First-order winners on some graphs with unit rates and a start of full
support (the second row allows any symmetric $Q$). The last column says where
each row is proved (Remark~\ref{rem:outcomes}).}
\label{tab:outcomes}
\end{table}

\begin{remark}[the outcomes table]\label{rem:outcomes}
Table~\ref{tab:outcomes} lists the outcomes for other graphs. Each row comes
from Theorem~\ref{thm:hull} applied to the exit rates of the connected faces.
So a row needs two things, namely the exit rates and the points that lie on
the lower-left boundary of the hull. The rows marked Lean are proved in Lean
(Section~\ref{sec:verify}). On $C_d$ two steps are by hand. A proper face that
is not an arc is disconnected, hence not admissible, and the growth of the
switch values follows from $2-2\cos x=x^2+O(x^4)$. The rows $P_d$ and
$\Theta_\varepsilon$ are proved in Appendix~\ref{app:examples}, for every
$d\ge3$ and for every $\varepsilon<\varepsilon_1$
(Propositions~\ref{prop:paths} and~\ref{prop:triangles}). The proof for
$\Theta_\varepsilon$ uses two resultants that we computed by computer algebra,
and it also shows that the bound $\varepsilon_1$ is sharp. In the row $P_4$ the
first sentence is the case $d=4$ of Proposition~\ref{prop:paths}, and the rest
is an exact computation of the law at the values of $N$ listed. A
well-separated community wins for a
long range of $\tau$, such as a triangle of $\Theta_\varepsilon$ from
$\tau\approx1$ to $\tau\approx9/\varepsilon$. The mode need not grow around its
first vertex. It leaves the leaf of the paw, and on the house and on $K_4$ plus
a vertex it jumps to a disjoint community. The second row explains the last
one. The state $5$ feeds $\{1,2,3,4\}$ unequally, so the point of
$\{1,2,3,4\}$ lies strictly left of the chord from $\{5\}$ to $[5]$, and the
winners cannot jump from $\{5\}$ straight to $[5]$. At
feasible $N$ the switch values are approached slowly
(Figure~\ref{fig:cycles}(b)), and on $P_4$ even the order of the phases
differs.
\end{remark}
